\documentclass[11pt]{article}
\usepackage[margin=1in]{geometry}
\usepackage{amsmath,amssymb,booktabs}
\usepackage[colorlinks=true,allcolors=blue]{hyperref}

\title{From Collision Histories to Computable Closure:\\
Proof-Guided Cumulants and Temporal Interaction Graphs}
\author{Working paper}
\date{\today}

\begin{document}
\maketitle

\begin{abstract}
Rigorous long-time kinetic limits propagate more than a smallness estimate for
correlations: they retain structured, time-oriented information about the
interaction histories that create those correlations. We formulate a
computational program that converts this proof mechanism into measurable
finite-scale closure diagnostics. The retained kinetic state is augmented by
projected connected cumulants and compressed temporal interaction graphs, while
the classical kinetic equation remains an explicit baseline. The principal
test is whether this augmented state predicts one-layer closure defect and
improves stable long-time rollout beyond models based only on the retained state
or unordered cumulant magnitudes. This document is a paper scaffold; its claims
remain hypotheses until the prescribed scaling, reversal, and ablation tests are
completed.
\end{abstract}

\section{Introduction}
State the computational closure problem, the arrow-of-time obstruction, and
the finite-scale research question.

\section{Mathematical blueprint}
Describe factorization, connected cumulants, time layering, partial expansion,
collision-history molecules, canonical layered gardens, and hydrodynamic
projection.

\section{Computational closure state}
For layer $\ell$, use
\[
S_\ell=(y_\ell,c_\ell,h_\ell,d_\ell,\sigma_\ell),
\]
where $y_\ell$ is the retained state, $c_\ell$ contains projected cumulants,
$h_\ell$ is a temporal history sketch, $d_\ell$ is closure defect, and
$\sigma_\ell$ quantifies uncertainty.

\section{Nonlinear waves to wave kinetics}
Specify the cubic NLS ensemble, four-wave kinetic baseline, cumulant estimators,
resonant hypergraph, and kinetic scaling sweep.

\section{Arrow-of-time tests and ablations}
Compare retained-state, cumulant-only, and oriented-history models on
forward/reverse pairs and history-shuffled controls.

\section{Hard-sphere transfer study}
Specify event-driven dynamics, weak marginal/cumulant estimators, collision
ancestry graphs, circuit rank, recollisions, and weak Boltzmann residuals.

\section{Results}
Report detectability, calibration, scaling, rollout stability, invariants,
uncertainty, and cost. This section is intentionally empty before experiments.

\section{Relation to the rigorous analysis}
Separate exact theorem objects from finite-scale estimators and graph proxies.
Explain why abstract cutting is used as a guide or certificate rather than as
the physical time integrator.

\section{Discussion}
Discuss negative results, application transfer, and the next kinetic-to-fluid
step.

\end{document}
